\section{The microstrip line: what a measurement of it can mean}
\label{sec:microstrip}

\noindent\emph{Status: draft -- the definitions, the lossless argument and the accounting identity below are
settled and are what the measurement rests on; the numbers they produce are reported separately and are not
quoted here.}

A microstrip is the simplest printed transmission line and the hardest to measure: an open, mixed-dielectric
guide whose ``mode'' is a convenience of the frequency range rather than an exact solution, driven at a port
whose impedance is not the impedance a catalogue quotes. This section fixes what is being measured before any
number is asserted. It sets out what the quasi-TEM assumption buys and costs, which of the several impedances
called $Z_0$ a transmission-line study must name, why moving the reference plane cannot repair a magnitude,
why a lossless solver can nonetheless return an attenuation, and how an attenuation is to be extracted from a
record that is contaminated by a standing wave. It closes with the acceptance criterion this measurement is
held to and with the identity that finally settles the question the estimators only bound.

\subsection{The quasi-TEM assumption and what it costs}
\label{sec:ms-quasithem}

A grounded dielectric slab supports a TM mode with no cutoff. Its lowest TM branch has a purely imaginary
transverse wavenumber at every frequency, so the transverse field of that mode solves Laplace's equation rather
than a wave equation, at zero frequency as much as at one gigahertz; the mode exists, and is quasi-static, at
arbitrarily low frequency. A strip printed on the slab inherits that property and loses its simplicity. The
fundamental mode of a strip on a grounded slab is a hybrid mode whose longitudinal components are not zero:
the strip edges carry a concentration of charge, because the tangential electric field must vanish on the
conductor and the field must therefore turn the corner there, so $E_z$ is non-zero on the edge and $H_z$ is
non-zero in the dielectric beneath it. Writing the mode as
\begin{equation}
  \mathbf{E}(x,y,z) = \left[\mathbf{e}_t(x,y) + e_z(x,y)\,\hat{\mathbf{z}}\right]e^{-\gamma z},
  \qquad \gamma = \alpha + j\beta,
  \label{eq:ms-hybrid}
\end{equation}
with $\mathbf{e}_t$ the transverse vector field and $e_z$ its longitudinal part, the quasi-TEM assumption is
the statement that $|e_z| \ll |\mathbf{e}_t|$, which holds when the transverse dimensions are small against the
guided wavelength and fails as the frequency rises. The assumption is therefore a statement about a range of
frequencies. It is not a statement about the physics of a grounded slab, which has no TEM mode at all: a TEM
mode requires two conductors and a single homogeneous dielectric between them, and the slab has one conductor
plane, one strip and two dielectrics.

What the assumption costs is a specific and non-transferable boundary-value problem. A quasi-static closed
form for a microstrip -- the Hammerstad--Jensen pair of expressions for $\eps_{\mathrm{eff}}$ and $Z_0$, or any
of its descendants -- is a fit to the numerical solution of Laplace's equation on a strip of zero thickness
lying on a laterally infinite substrate of one permittivity with free space above, in a domain that radiates
nothing because it is not bounded at all. The box in this solver is nothing of the kind. It is closed on every
face, and the only faces that are not perfect conductors are the matched layers that terminate the line at
both ends and that line the lateral and upper boundaries. What the solver launches is therefore the
fundamental mode of a shielded, lossy-walled box that contains a strip: it has a discrete spectrum, its
transverse profile is influenced by walls a finite distance away, and it is the object that must be measured,
because it is the object that propagates. Hammerstad--Jensen is consequently used here as a diagnostic and not
as an oracle. It answers whether the geometry, the effective permittivity scale and the impedance are of the
right order, and whether a port-plane measurement is in the right units; it does not answer whether the mode
the solver propagates is the open-line mode, because it solves a different problem. Disagreement with it is
evidence about the box, not automatically an error in the solver, and agreement with it is not evidence that
anything beyond the geometry is right.

\subsection{Three impedances, and which one is meant}
\label{sec:ms-z0}

Three quantities are routinely called the characteristic impedance of the same line, and they are equal only
under an assumption that a real measurement is testing. The first is the ratio of the line voltage to the line
current at a plane,
\begin{equation}
  Z^{VI} = \frac{V}{I}, \qquad
  V = -\int_{\text{strip}}^{\text{ground}} \mathbf{E}\cdot \dd\mathbf{l}, \qquad
  I = \oint_{\text{strip}} \mathbf{H}\cdot \dd\mathbf{l},
  \label{eq:ms-zdefs}
\end{equation}
with $V$ in volts, $I$ in amperes and $Z^{VI}$ in ohms. The second and third are built on the power $P$ in
watts crossing the plane, $Z^{PI} = 2P/|I|^2$ and $Z^{PV} = |V|^2/2P$. For a single wave travelling in one
direction they coincide exactly, since $P = \tfrac12\operatorname{Re}(VI^{*})$ and $V = Z_0 I$ give
$Z^{VI} = Z^{PI} = Z^{PV} = Z_0$, the real characteristic impedance of the mode. With a reflected wave of
coefficient $\Gamma$ present at the plane the three separate,
\begin{equation}
  Z^{VI} = Z_0\,\frac{1+\Gamma}{1-\Gamma}, \qquad
  Z^{PI} = Z_0\,\frac{1-|\Gamma|^2}{|1-\Gamma|^2},
  \label{eq:ms-zgamma}
\end{equation}
and the spread between them is not a nuisance to be averaged away but a measurement: it is the non-modal
content of the field at that plane, reflection first and higher modes with it. A study that reports one number
called $Z_0$ without saying which definition it used has reported a quantity whose disagreement with a model
cannot be interpreted.

The third quantity in play is the one a user's S-parameters are actually quoted on: the lumped impedance of
the port plane, $Z_{\mathrm{plane}} = V/I$ evaluated where the source sits. It is complex, and its imaginary
part is the reactance of the lumped source and of the near field it drives, not a property of the line. Writing
$Z_{\mathrm{plane}} = R + jX = |Z_{\mathrm{plane}}|e^{j\phi}$ factorises the complex number into a modulus and a
real part,
\begin{equation}
  |Z_{\mathrm{plane}}| = \frac{R}{\cos\phi}, \qquad \phi = \arctan\frac{X}{R},
  \label{eq:ms-zplane}
\end{equation}
so the modulus is inflated by the source's own reactance and is shifted off the resistive value by the same
angle. The modulus is therefore not the characteristic impedance, it is not the resistive load the line
presents, and it must never be printed as though it were: the quantity to compare against a model is
$\operatorname{Re} Z_{\mathrm{plane}}$, read at the same plane and with the same definition, and the phase is
reported beside it as what it is -- a measure of how reactive the launch is, and hence of how far the port is
from the pure travelling wave that \eqref{eq:ms-zgamma} needs.

\subsection{De-embedding and the reference plane}
\label{sec:ms-deembed}

Moving a scattering parameter along a line of known propagation constant is an exact operation,
\begin{equation}
  S_{11}(z+d) = S_{11}(z)\,e^{-2\gamma d}, \qquad
  S_{21}(z+d) = S_{21}(z)\,e^{-\gamma d}\ \text{at each end},
  \label{eq:ms-deembed}
\end{equation}
with $d$ the shift in metres and $\gamma$ as in \eqref{eq:ms-hybrid}. Because $\gamma$ is complex, the shift
scales the magnitude by $e^{-2\alpha d}$ and leaves it alone only when $\alpha = 0$. In a lossless model
$\gamma = j\beta$ is purely imaginary, and then \eqref{eq:ms-deembed} changes the phase of every scattering
parameter and nothing else. Two consequences follow, and both are used later. First, $|S_{11}|$ is identical at
every reference plane of a lossless line, and so is $|S_{21}|$; a discrepancy in magnitude between a
port-plane measurement and a model therefore cannot be relocated by de-embedding, and any attempt to explain
one that way is an attempt to smuggle an attenuation into a bookkeeping step. Second, because the only
magnitude-carrying part of the shift is the $\alpha$ under test, using de-embedding to close a magnitude gap
would assume the answer. What de-embedding does own is the phase, and it owns it completely: reading a field
quantity half a cell from where it is defined appears in \eqref{eq:ms-deembed} as a linear phase, which is why
the phase of a port-plane impedance and the position of a station must be fixed before either is compared
with anything.

\subsection{Where the damping can come from in a lossless model}
\label{sec:ms-lossless}

The time-domain solver contains no mechanism for ohmic or dielectric loss, and the claim is checkable
mechanism by mechanism rather than by assertion. The update is a plain leapfrog pair, and the only step in it
that is not unitary is the multiplication of $\mathbf{e}$ and $\mathbf{h}$ by the matched-layer factors of
\eqref{eq:pml-factor} immediately after their own updates. Those factors are $\exp(-\sigma_e\Delta t/\eps_0)$
and $\exp(-\sigma_m\Delta t/\mu_0)$ with $\sigma_m = \sigma_e\mu_0/\eps_0$, so they are equal wherever the
profile is non-zero, they lie in $(0,1]$, and they are exactly one wherever $\sigma_e$ vanishes -- and
$\sigma_e$ is non-zero only inside the layers, so every multiplier applied in the physical interior is
identically one. The material matrices that carry the substrate are built from a real relative permittivity
and a real relative permeability and an interface rule for averaging across a material jump; the routine that
assembles them has no conductivity argument at all, and the conductivity that exists elsewhere in the package
belongs to the frequency-domain resonance helper, which no time-domain study calls. The strip is an ideal
conductor: the edges that lie in it are flagged and the tangential electric field on them is set to zero after
every update, with no surface impedance, no skin depth and no metal thickness anywhere in the model, so the
conductor cannot dissipate by construction. Finally the discrete energy identity of
Section~\ref{sec:energy} holds to round-off over a full run, which means the interior update moves energy
between degrees of freedom and destroys none of it.

The consequence is a sharp statement about what a decay can mean. Since no operator in the interior lowers the
energy, the amplitude of the launched mode can fall only by transferring power out of it, and the only place in
the box where power can finally leave is the matched layers; every other face is a perfect conductor. A decay
of the guided mode with distance along the line is therefore a boundary sink, and the only candidate sinks are
the matched layers. This does not yet say what the decay rate is, and it does not rule out a second
possibility: power transferred out of the launched mode into another mode of the same closed box stays in the
box, bounces between perfect conductors, and returns. In that case the station record does not decay at all,
it oscillates, and the estimators of the next subsection are what turn the oscillation into a spurious
attenuation. The two possibilities are distinguished by measurement, not by argument.

\subsection{Five estimators of one attenuation}
\label{sec:ms-estimators}

The same record -- the complex line voltage at a set of station planes spaced along the line -- supports five
distinct estimators of the attenuation, and they do not agree: on the default case they differ by a factor of
order two. The first is a one-wave joint fit, which models every station phasor as one forward exponential,
\begin{equation}
  V_n = A\,e^{-\gamma z_n}, \qquad n = 1,\dots,N,
  \label{eq:ms-onewave}
\end{equation}
and returns $\alpha$ and $\beta$ together by least squares, with $A$ in volts, $z_n$ in metres. The second is
a two-wave joint fit, which adds the backward wave and so allows the record to contain a standing wave,
\begin{equation}
  V_n = A\,e^{-\gamma z_n} + B\,e^{+j\beta z_n},
  \label{eq:ms-twowave}
\end{equation}
at the price of a non-linear search in $\beta$. The third is the endpoint ratio, $|S_{21}| = |V_N|/|V_1|$,
which assumes one travelling wave and nothing else and is therefore exact when \eqref{eq:ms-onewave} holds and
wrong in proportion to the reflected content when it does not. The fourth is the forward-subtracted linear
fit: fit \eqref{eq:ms-twowave}, subtract the fitted backward wave, and then take an ordinary least-squares line
through $\ln|V_n|$ against $z_n$ for the remaining forward amplitude. The fifth is the model-free end ratio of
$|V|$, the ratio of the raw amplitudes at the two extreme planes of the record, with no fit, no phase model and
no removed wave at all -- the number a laboratory would read off an instrument.

That these five disagree is the point. Each is unbiased under a different assumption about what else is in the
record, and a fit is not more trustworthy than a ratio merely because it uses more points. The one-wave fit is
the estimator most exposed to the standing wave, because a single exponential has no way to represent a
ripple: it must absorb one, and the only free parameter left to absorb it is the slope. Writing the record as
a forward wave plus a reflection, $\ln|V(z)| = \ln|A| - \alpha z + \ln|1+\Gamma e^{2j\beta z}|$, the last
term is periodic in $z$ with the ripple period, and over a window it contributes a bias whose size is set by
the ripple amplitude and by how much of a period the window covers, not by the quality of the fitting. A short
span, or a span that happens to cover a fraction of a period, biases the slope; no amount of care in the
solver or in the fit removes it. This is why the disagreement is a property of the record and not of the
estimator's competence, and why none of the five can be promoted to the answer on the strength of its own
output. What is needed is an estimator that has no free parameter into which a ripple can vanish.

\subsection{The acceptance criterion and its resolution}
\label{sec:ms-criterion}

The criterion this measurement was originally held to was an absolute one: $|S_{21}|$ within $1$ per cent of
unity, that is within the interval $0.990$ to $1.010$. An absolute target is only a test if the measurement can
resolve it, and this one cannot. The argument is not that the target is hard to hit; it is that the estimators
available on this record disagree with each other by nearly twice the width of the target, and that the
disagreement has a cause no estimator on this record can remove.

The cause is the port reflection. At $|\Gamma| = 0.0060$ in amplitude a standing wave moves the raw envelope of
the field by $\pm 1.2$ per cent, that is by $2.4$ per cent peak to peak, at the ripple period. Every model-free
attenuation on this record is a ratio of the field at two stations, and each such ratio inherits that movement:
the raw envelope falls $3.0$ per cent across the $9.60$ mm of measured line, so a single ratio is not good to one
per cent, while the half-span fits, which average over several ripples, are individually good to
$1.4\times 10^{-3}$ in $\ln|V|$. That is why the disagreement does not shrink when the fit is done more
carefully: the two halves of the record agree with each other and disagree with the ratios.

The floor of \eqref{eq:ms-criterion-number} is worth stating exactly, because it is easy to misread as a
property of the box of Section~\ref{sec:msres-box}. It is not. The clearance and lid sweeps of the earlier study
put a floor of
\begin{equation}
  \alpha_{\mathrm{dB}} = 0.0139 \pm 0.0026\ \text{dB/mm}
  \label{eq:ms-criterion-number}
\end{equation}
over a $14.4$ mm span, where it corresponds to $|S_{21}| = 0.977$. It is a bound on the \emph{fitted} attenuation
of those sweep boxes over that span, and a rate obtained from audited wall work is a different estimator, so the
two must never be compared. The reference box is a different geometry over
a different span, and its fitted attenuation is $0.02145$ dB/mm, not $0.0139$.

The decisive statement is the spread among the estimators on the reference box itself. The five rows of
Table~\ref{tab:ms-alpha} span $0.01160$ to $0.02164$ dB/mm, a factor of $1.86$, and the span is not noise about
a common value: the two model-free rows sit together at $0.01160$ and $0.01183$, the two one-wave fits sit
together at $0.02145$ and $0.02164$, and the gap between the two groups is exactly the standing-wave bias
analysed in Section~\ref{sec:msres-alpha}. Charged over the $21.6$ mm of line between the port plane and the far
plane, the span runs from
\begin{equation}
  |S_{21}| = 10^{-0.01160 \times 21.6/20} = 0.972
  \quad\text{to}\quad
  |S_{21}| = 10^{-0.02164 \times 21.6/20} = 0.948 .
  \label{eq:ms-criterion-span}
\end{equation}
The interval a $1$ per cent target asks for, $0.990$ to $1.010$, lies entirely above that span, and no estimator
on this record places $|S_{21}|$ inside it. The resolution of the record is therefore $2.5$ to $3$ per cent in
$|S_{21}|$, and a target stated to one per cent would be graded by which estimator was chosen rather than by
whether the solver is right.

The criterion is therefore revised rather than waived, and the revision is stated as three requirements that
must all hold. First, $|S_{21}|$ must be consistent with a lossless \emph{material} model within the stated
measurement uncertainty, which is the estimator span of \eqref{eq:ms-criterion-span} and is a statement about
what the record can resolve, not a permission to be wrong by any amount. Second, the attenuation that remains
must be attributed to its dominant mechanism, with the attribution accounting for the magnitude and not merely
for the sign. Third, the mechanism must be identified, by name and by the part of the box responsible, so that
the shortfall from unity is a known numerical effect rather than an unexplained residual; and where the record
cannot settle which of several physical mechanisms it is, the criterion requires that the candidates be named
and the measurement that would separate them be stated. A criterion in that form can fail: it fails if the
measured $|S_{21}|$ leaves the span, or if no mechanism accounts for the number, or if the mechanism named is
one the solver cannot produce, or if the candidates left open are not bounded. What it refuses to do is grade a
solver against a tolerance narrower than the spread of its own measurement.

\subsection{The accounting identity}
\label{sec:ms-accounting}

The estimator that has no free parameter is an energy balance, and it is available here only because the model
is lossless. In a lossless box the discrete energy identity of Section~\ref{sec:energy} is exact, so over any
window in which the field has reached a steady state every joule injected has left through a boundary, and the
work removed by a boundary is not an estimate but an identity. For a wall family $i$ that removes work $W_i$
(joules) during a window in which the run injects $U_{\mathrm{in}}$ (joules), and that runs along $L_i$ metres
of the line, the per-unit-length rate is
\begin{equation}
  \alpha_i = \frac{W_i}{2\,U_{\mathrm{in}}\,L_i},
  \label{eq:ms-accounting}
\end{equation}
with $\alpha_i$ in nepers per metre. The factor of two is the definition of an amplitude attenuation: a wave
carrying power $P$ (watts) loses $2\alpha P$ per metre, so over the window
$W_i = 2\alpha_i L_i \int P\,\dd t = 2\alpha_i L_i U_{\mathrm{in}}$, which is \eqref{eq:ms-accounting}
rearranged. For reference, one neper per metre is $8.6859$ dB per metre, or $8.6859\times10^{-3}$ dB per
millimetre.

Three properties of \eqref{eq:ms-accounting} matter, and they are the reason it outranks the fits of
Section~\ref{sec:ms-estimators}. It has no free parameter, so a standing-wave residual cannot be absorbed into
it: a ripple moves energy back and forth inside the box and nets to nothing in $W_i$ over a full period, which
is exactly the behaviour the slope of a fit cannot deliver. It makes no model of the mode, so it does not
assume the field is quasi-TEM, single-mode, or anything else -- it counts joules, and joules are conserved by
the scheme to round-off. And it is local in the sense that matters: it says which family of walls removed the
power, so the attribution required by Section~\ref{sec:ms-criterion} is produced by the same computation that
produces the rate. A rate from an identity is a stronger measurement than a rate from a fit, and the
difference is not a matter of precision but of what each can represent.

The identity has limits, and they should be stated with it. It needs a steady state, so the warm-up steps
before the field has settled must be excluded from the window, and the injected power must be the power the
box actually retains rather than the power the source nominally delivers. It returns one global rate for the
run, not a rate per station, so it cannot show where along the line the loss accumulates. And it is a
statement about boundaries, so it measures how much power left through which wall, not why the mode was not
confined in the first place: the normalisation by $L_i$ is the only step that converts a boundary flux into a
per-unit-length rate, and it is valid only for that part of the wall that the line actually runs alongside.
Charging a wall that spans the whole box to a shorter measurement span gives an upper bound on the rate and
must be reported as one. Within those limits, the accounting converts the question from ``what does the slope
of this record say'' to ``where did the power go'', and the second question has an answer that does not depend
on how the record was fitted.
